\documentclass[11pt,a4paper]{article}
\usepackage{fontspec,amsmath,amssymb,amsthm,microtype}
\setmainfont{lmroman10-regular.otf}[BoldFont=lmroman10-bold.otf,ItalicFont=lmroman10-italic.otf,BoldItalicFont=lmroman10-bolditalic.otf]
\usepackage[margin=25mm]{geometry}
\usepackage[unicode,hidelinks]{hyperref}
\hypersetup{pdfauthor={LangevinLean Project}}
\setlength{\emergencystretch}{3em}
\setlength{\parskip}{0.4em}
\setlength{\parindent}{0pt}
\allowdisplaybreaks[2]
\raggedbottom
\title{\texorpdfstring{Physical-\allowbreak{}time particle integral solutions and synchronous contraction}{Physical-time particle integral solutions and synchronous contraction}}
\author{LangevinLean Project}
\date{}
\begin{document}\maketitle
Draft mathematical exposition; these notes may contain mistakes.


The particle-average potential has the original pair curvature and smoothness constants, so the physical-time kinetic system can reuse the shared existence and nonlinear contraction theorems directly. The temperature changes the additive noise amplitude; the synchronous contraction does not depend on that amplitude because the common forcing cancels. The results here are explicit path statements, not stationarity statements.



The average force meets the shared kinetic existence and contraction hypotheses. Common-input cancellation explains why physical-time deterministic contraction does not deteriorate with particle number or total state dimension.



Earlier results used here: \cite[Section~3.1]{LangevinLean2026Kinetic}; \cite[Section~3.6]{LangevinLean2026Kinetic}.

\section{\texorpdfstring{Verify differentiability and smoothness of the average}{Verify differentiability and smoothness of the average}}


The lifted potential is differentiable whenever the pair potential is differentiable. Multiplying its value by the constant \(1/N\) preserves differentiability, yielding differentiability of the average potential. The gradient has already been identified pointwise with the particle force. Equality at every position identifies the two vector fields. The average-gradient smoothness bound therefore gives the same Lipschitz estimate for the particle force. The result uses the original pair constant \(L\), independent of \(N\).



The gradient identity is essential for applying a potential-based contraction theorem: it ensures the defined force is the derivative of the stated average potential. Smoothness is then inherited with the original pair constant, while strong convexity is inherited by averaging the lifted inequality.



Earlier results used here: \cite[Section~3.1.1]{LangevinLean2026IPLA}; \cite[Section~3.1.3]{LangevinLean2026IPLA}.

\section{\texorpdfstring{Apply the kinetic integral-\allowbreak{}equation construction to the real particle force}{Apply the kinetic integral-equation construction to the real particle force}}


With \(N>0\), the force is globally Lipschitz. The shared kinetic existence theorem therefore constructs a continuous phase path satisfying \(Y_t-Y_0=-\int_0^tD_F(Y_s)ds+(B_t,0)\) on any finite nonnegative horizon, for every continuous velocity forcing starting at zero. The integral equation uses inverse mass one because the force already equals \(\nabla\bar U_N\). Choosing physical-time Brownian amplitude \(\sqrt{2\gamma/N}\) gives the intended scaled particle forcing, but a complete probabilistic strong-solution result also identifies the driving Brownian law and independence conditions.



The forcing is any continuous velocity input in the deterministic existence statement. The specified Brownian choice supplies the desired thermal process afterward. This separates solving the integral equation from identifying the probability law of its driving source.



Earlier results used here: \cite[Section~3.1.1]{LangevinLean2026IPLA}; \cite[Section~3.1.1]{LangevinLean2026Kinetic}; \cite[Section~3.1.3]{LangevinLean2026Kinetic}.

\section{\texorpdfstring{The nonlinear contraction rate is independent of particle dimension}{The nonlinear contraction rate is independent of particle dimension}}


Apply the shared potential contraction theorem to \(\bar U_N\). Its required hypotheses are derived by the average-potential differentiability, strong-convexity and smoothness theorems. Rewrite each gradient in the difference equation using \(\nabla\bar U_N=F\). With equal initial velocities, \(\gamma>0\), and \(\gamma^2\ge m+L\), the resulting bound is \(\|x_t-y_t\|\le\sqrt2e^{-mt/\gamma}\|x_0-y_0\|\). Neither \(N\) nor \(d_\theta+Nd_x\) appears in its rate. The theorem still explicitly requires continuous paths satisfying the noise-cancelled derivative equations; it does not replace construction of a stationary kinetic reference or taking the Wasserstein infimum.



The common additive noise cancels even though its amplitude depends on particle count. The remaining difference equation sees only the average-potential geometry. The dimension and temperature therefore affect equilibrium and discretization moments, while leaving the synchronous physical-time rate unchanged.



Earlier results used here: \cite[Section~3.1.1]{LangevinLean2026IPLA}; \cite[Section~3.6.5]{LangevinLean2026Kinetic}.

\begin{thebibliography}{ULA-LSI}
\bibitem[IPLA-SC]{LangevinLean2026IPLA}
LangevinLean Project. \emph{Interacting Particle Langevin under Strong Convexity}. LL-03, 2026. Draft report, content revision c8b5d67ceb43. \url{https://odakyildiz.github.io/LangevinLean/docs/pdf/IPLA/FullProof.pdf}.
\bibitem[Kinetic]{LangevinLean2026Kinetic}
LangevinLean Project. \emph{Kinetic Langevin Monte Carlo}. LL-06, 2026. Draft report, content revision 57bac95a396b. \url{https://odakyildiz.github.io/LangevinLean/docs/pdf/KineticLangevin/FullProof.pdf}.
\end{thebibliography}

\end{document}
